\subsection{Beschaffung der Testdaten}\label{subsec:beschaffung_der_testdaten}
Die Testdaten wurden aus verschiedenen Quellen beschafft. Die Geodaten (GeoTIFF-Dateien) wurden von verschiedenen Quellen beschafft. Für die Alpen wurden die Daten des Copernicus Programms der EU \todo{abkürzung} bezogen. Diese Entscheidung wurde Aufgrund der Ländergrenze zwischen Deutschland und Österreich, welche sich in dem Kartenausschnitt befindet, getätigt. Als Kartenmodell wurde hier die \enquote{Digital Elevation Model (DEM) for Europe at 30 meter} gewählt. Aus ihr können Kartenausschnitte von Europa im mit einer Auflösung von 30m pro Pixel heruntergeladen werden. Um die Performance des Tests zu verbessern wurde der Kartenausschnitt auf eine Pixelgröße von ca. 90m umgerechent. Die Kartendaten für den Südschwarzwald wurden auch aus dieser Quelle gewonnen. Das liegt daran, dass auch Teile der nördlichen Schweiz bei Schaffhausen abgebildet werden. Für den Südschwazwaldtest wurde die Kartengenauigkeit nich verringert, sondern auf 30m belassen.

Für den Feldberg wurden die Daten von der Landesvermessung Baden-Württemberg bezogen. Der Grund hierfür ist, dass die Daten dort genauer sind. Alle Daten wurden in einer Auflösung von 0.25m pro Pixel bezogen. Da diese Auflösung zu hoch für die Testzwecke ist, wurden die Daten auf eine Auflösung von 10m pro Pixel heruntergerechnet.

Die Daten über die realen Gipfel wurden aus verschiedenen Quellen bezogen. Da in den Daten nicht flächendeckend  die Prominenz und Dominaz definiert wruden, mussten bei der Erstellung der Realdatensätze verschiedene API-Quellen und manuelle Suche kombiniert werden. Zum einen wurden die Daten von der OpenStreetMap API bezogen. Diese Daten wurden mit Hilfe der Overpass API abgefragt. Es wurden alle Punkte mit dem Tag \textit{natural=peak} innerhalb der Kartenausschnitte abgefragt. Diese Daten wurden wurden zunächst mit den Daten von Wikidata abgeglichen und durch diese ergänzt. Wikidata bietet eine API an, über die für bestimmte Koordinaten die Gipfel abgefragt werden können. Mit diesem Abgleich konnten über mehrere Kartenbereichen in den Alpen im Durchschnitt 80\% der Gipfel mit Prominenz über 100m gefunden werden. Dies leigt nicht daran, dass die Gipfel selbst unbekannt sind, sondern daran, dass die Prominenz und Dominanz nicht in den Datenquellen definiert wurden. Daher wurden die Daten von Hand ergänzt und auf Vollständigkeit überprüft. Zum besseren Analyseren der Tests wurden die Gipfel, die nicht den Kriterien für einen Gipfel entsprechen, als \textit{Fake Peaks} definiert und in die Testdaten aufgenommen. Es wurden Punkte ausgewählt, die in der Nähe von realen Gipfeln liegen, aber nicht die Kriterien für einen Gipfel erfüllen oder von den APIs als Gipfel, die den Bedingungen nicht entprechen erkannt wurden. Diese Daten helfen bei dem schnellen finden von Fehlern des Algorithmus, da sie zeigen, welchen Gipfel der Algorithmus fälschlicherweise als Gipfel identifiziert hat und nicht nur seine Koordinaten.
Um die Daten manuell zu ergänzen und zu überprüfen muss die Karte von Hand durchsucht werden. Dabei stellt sich die Frage, welche Karte verwendet werden sollte, um die genausten testdaten zu erhalten. Dabei spielen Faktoren wie die allgemeine Genauigkeit der Höhenlinien und Anzahl der makierten Gipfel eine Rolle. Die Höhenlinien sind wichtig um die Prominenz abzuschätzen, wenn keine Realwerte dafür existieren. Bei wie vielen Gipfel keine Realwerte in den Datenquellen existieren hängt von vielen Faktoren ab. Dabei gilt meist, dass desto berühmter, höher die Gipfel und desto höher deren Prominenz und Dominanz ist umso wahrscheinlicher existieren Realwerte für diese Berge und umso besser sind diese. Ersteres spielt nach der Erfahrung eine größere Rolle. So ist beispielsweise die Höhe der Zugspitze in der Wikidata Datenbank auf Centimeter genau definiert, während die Höhe der meisten anderen Gipfel nur auf Meter genau definiert ist. Es gibt aber auch Fälle, wie beispielsweise den Daniel, bei der die Höhe nur auf 10 Meter genau definiert ist, obwohl sie ein eine große Prominenz von über 500m besitzt.

Es wurden verschiedene Karten verglichen, wie die OpenStreetMap Karte, die Karte von Google Maps und die Karte von Bing Maps. Es wurde festgestellt, dass die OpenStreetMap Karte die genausten Höhenlinien und die meisten markierten Gipfel enthält. Daher wurde sie als Hauptquelle für die manuelle Ergänzung der Testdaten verwendet. Die Werte für die Prominenz und Dominanz wurden dabei durch Abschätzen der Höhenlinien auf der Karte ergänzt, wenn nicht schon durch die APIs gegeben.